\section{Compact-source sheaf categories and homotopy invariance}
\label{sec:categories}

The sheaves used below can have very complicated stalks.  Their finiteness
comes from their sources: they are built, by finitely many operations, from
proper images of constant sheaves on finite polyhedra.  We first establish
duality on exactly these categories, and then prove the cylinder relation
that will make their Witt groups homotopy invariant.

\subsection{Spaces, coefficients, and generators}

All sheaves in this section have coefficients in the field $\R$, regarded
algebraically.  A \emph{test space} is a locally compact, finite-dimensional,
separable metrizable space of the following kinds: a finite polyhedron, a
space homeomorphic to one, an open subset of such a space, or a finite
product of these spaces.  Points and the empty space are included.  The
spaces to which Verdier duality is applied will always be test spaces.

Write
\[
 \mathsf D(X)=D^b(\operatorname{Sh}_{\R}(X)).
\]
For a complex $A$, the assertion that its support is contained in a compact
set $K\subset X$ means that all its cohomology sheaves vanish on $X\setminus
K$.  We use this formulation even when the set of points with a nonzero
cohomology stalk is not closed.  All tensor products are derived tensor
products over $\R$; tensoring sheaves is exact over this field.

\begin{definition}[The test categories]
\label{def:test-categories}
Let $\Tzero(X)$ be the smallest full replete triangulated subcategory of
$\mathsf D(X)$ containing
\begin{equation}
 \label{cat:generators}
 Rf_*\const{\R}{B},
 \qquad f:B\longrightarrow X,
 \qquad B\text{ a compact finite polyhedron},
\end{equation}
where $f$ is any continuous map.  Let $\T(X)$ be its closure under direct
summands in $\mathsf D(X)$.

When $X$ has a specified semialgebraic model, define $\Pzero(X)$ and
$\Pcat(X)$ in the same manner, requiring $B$ to be a compact polyhedron in
a Euclidean space and $f$ to be continuous semialgebraic in the specified
models.  For an arbitrary open subset $U$ of a semialgebraic model, the
condition on $f:B\to U$ is that its composite into that model be
semialgebraic.  The open subset itself need not be semialgebraic.

An allowed map for $\T$ is a continuous map of test spaces.  For $\Pcat$ we
use continuous semialgebraic maps in the models, their restrictions to the
open subsets under consideration, and inclusions of such open subsets.
The condition on a map used in the arguments below is
that composition carry every allowed compact-source generator to an
allowed compact-source generator.  Homotopies for $\Pcat$ satisfy this
condition in their two variables; the continuous semialgebraic homotopies
used below do so.
\end{definition}

The sheaf operations used below satisfy their finite-dimensionality
hypotheses for every allowed map, including a nonproper map.  To see
this explicitly, every test space is homeomorphic to a locally closed
subset of some $\R^N$: realize its finite polyhedral factors as closed
Euclidean subsets, take the prescribed relative open subsets, and
then their finite product.  Fix such an embedding of the source $X$
of a continuous map $g:X\to Y$.  Since $Y$ is Hausdorff, each fiber
$Z=g^{-1}(y)$ is closed in $X$, hence locally closed in the same $\R^N$.
For its embedding $i:Z\hookrightarrow\R^N$, extension $i_!$ is exact:
factor $i$ as a closed inclusion followed by an open inclusion.
For every sheaf $F$ on $Z$, composition of derived proper images gives
\[
 R\Gamma_c(Z;F)\simeq R\Gamma_c(\R^N;i_!F).
\]
The right side has no cohomology above $N$, by
\cite[Lemma~11.7.1]{SchapiraNotes}.  This bound is independent of $y$.
Thus every such $g$ has finite $c$-soft dimension, by
\cite[Definition~11.3.1 and Remark~11.3.3]{SchapiraNotes}.
The composition formula used here is established before that
finite-dimension hypothesis is imposed
\cite[Equation~(11.3.1)]{SchapiraNotes}, so this verification does
not assume the duality formalism it is used to justify.

In particular the maps in \eqref{cat:generators} are proper, because
their sources are compact and their targets are Hausdorff, and their
derived images are bounded.  Proper base change identifies their
stalk cohomology with the cohomology of compact fibers.  Boundedness
does not imply that these stalks are finite dimensional.
For the general locally compact sheaf formalism see
\cite[Chapter~III]{KashiwaraSchapira1990} and
\cite[\S\S11.3--11.4]{SchapiraNotes}.

We work with essentially small categories.  This causes no size issue:
there is a set of finite polyhedra up to isomorphism and of continuous maps
from each into a fixed $X$; the derived category is locally small.  Finite
constructions using their morphisms, followed by the choice of summands,
still have a set of isomorphism classes.  Choose representatives and then
take replete subcategories when convenient.

\begin{lemma}[Simplex generators and summands]
\label{cat:simplex-generators}
In either definition, closed simplices suffice in place of all compact
finite polyhedra.  Every object of $\T(X)$, respectively $\Pcat(X)$, is a
direct summand of an object of $\Tzero(X)$, respectively $\Pzero(X)$.
In particular these summand closures are triangulated and thick.
\end{lemma}

\begin{proof}
Fix a finite triangulation of $B$.  For a closed simplex $\sigma$, let
$j_\sigma:\sigma^\circ\hookrightarrow\sigma$ and
$i_\sigma:\partial\sigma\hookrightarrow\sigma$ be the inclusions.  The
open-complement triangle is
\begin{equation}
 \label{cat:simplex-boundary-triangle}
 (j_\sigma)_!\const{\R}{\sigma^\circ}
 \longrightarrow\const{\R}{\sigma}
 \longrightarrow(i_\sigma)_*\const{\R}{\partial\sigma}
 \longrightarrow
 (j_\sigma)_!\const{\R}{\sigma^\circ}[1].
\end{equation}
Induction on dimension expresses the first term using constants on closed
simplices: apply the same argument to the finite triangulation of the
boundary, starting with points.  Filter $B$ by its skeleta.  On each
skeleton the complement of the previous skeleton is a finite disjoint
union of open simplices.  Its constant sheaf extended by zero is the
finite sum of the first terms of
\eqref{cat:simplex-boundary-triangle}, extended into that skeleton.
The open-complement triangles for this filtration therefore build
$\const{\R}{B}$ from constants on the closed simplices of $B$.
Pushing these finite triangles along $f$ proves the assertion.  For the
semialgebraic categories all restrictions and simplex inclusions just
used are semialgebraic.

For the second assertion, the class of summands of objects in a full
triangulated subcategory is itself triangulated.  Indeed, suppose
$A\oplus A'=C$ and $B\oplus B'=D$ belong to that subcategory, and let
$u:A\to B$.  Fullness places the block map $u\oplus0:C\to D$ in the
subcategory.  A cone on this map is
\[
 \Cone(u)\oplus B'\oplus A'[1].
\]
Hence $\Cone(u)$ is a summand of an object of the subcategory.  Shifts and
finite sums have the same property, and an iterated summand is a summand.
This proves the claim and also proves thickness.
\end{proof}

\subsection{Duality on the generated categories}

Let $a_X:X\to\pt$.  We use the dualizing complex and Verdier operation
\[
 \omega_X=a_X^!\R,
 \qquad
 \D_X A=R\mathcal Hom(A,\omega_X).
\]
The operation is initially understood in the ambient derived category.
Its boundedness and the invertibility of its biduality map on our
categories will be proved, not assumed.

We recall precisely the pieces of the sheaf formalism used here.  For
locally compact spaces and maps of finite $c$-soft dimension, $Rf_!$ has
right adjoint $f^!$, composition identifies $f^!\omega_Y$ with
$\omega_X$, and the counit gives the trace
\begin{equation}
 \label{cat:trace}
 \operatorname{Tr}_f:Rf_!\omega_X\longrightarrow\omega_Y.
\end{equation}
The internal duality isomorphism is
\begin{equation}
 \label{cat:verdier-formula}
 Rf_*\D_X A\xrightarrow{\sim}\D_Y Rf_!A
 \qquad(A\in\mathsf D(X)).
\end{equation}
It is adjoint to evaluation followed by the trace, and is compatible with
the evaluation map $\eta_A:A\to\D_X^2A$ whenever the expressions are
bounded.  For an open inclusion $j$, $j^!=j^{-1}$, so duality commutes
with open restriction.  These are the adjunction and internal duality
statements in \cite[Theorem~11.4.2, Propositions~11.4.3--11.4.4 and
11.4.6(b)]{SchapiraNotes}.  Proper base change, the projection formula,
and compact-support K\"unneth will also be used in their natural forms
\cite[Theorems~11.3.5--11.3.6 and Corollary~11.3.11]{SchapiraNotes}.
None of these statements asserts reflexivity of every bounded sheaf
complex.

\begin{proposition}[Duality and compact support]
\label{prop:test-duality}
The Verdier operation restricts to a triangulated duality on each of
$\Tzero(X)$, $\T(X)$, $\Pzero(X)$, and $\Pcat(X)$.  In particular, for
every object $A$ in one of these categories, $\D_X A$ belongs to the same
category and the canonical map
\[
 \eta_A:A\xrightarrow{\sim}\D_X^2A
\]
is an isomorphism.  There is a compact $K\subset X$ outside which both
$A$ and $\D_X A$ vanish.
\end{proposition}

\begin{proof}
We first calculate on a closed $n$-simplex $\Delta$, choosing an
orientation of its interior.  With $j:\Delta^\circ\hookrightarrow
\Delta$,
\begin{equation}
 \label{cat:simplex-duality}
 \D_\Delta\const{\R}{\Delta}
   =\omega_\Delta
   \simeq j_!\const{\R}{\Delta^\circ}[n],
 \qquad
 \D_\Delta\bigl(j_!\const{\R}{\Delta^\circ}[n]\bigr)
   \simeq\const{\R}{\Delta}.
\end{equation}
Here is the local calculation, including the boundary.  At an interior
point the compactly supported cohomology of a sufficiently small open
ball is $\R$ in degree $n$ and zero otherwise.  At a boundary point a
cofinal system of relative neighborhoods is homeomorphic to a half-ball
with its flat face retained.  All its compactly supported cohomology
vanishes.  For example, its one-point compactification is a closed ball
and has zero reduced cohomology; this also applies at corners after a
local homeomorphism.  The local description of the dualizing complex
therefore gives the first isomorphism in
\eqref{cat:simplex-duality}.  The orientation identifies its interior
local system with the constant sheaf.  For the dualizing-complex
formulation of orientation on a manifold without boundary, see
\cite[\S11.7]{SchapiraNotes}; the relative half-ball calculation above
establishes the boundary case needed here.

For the second isomorphism, open-image duality gives
\[
 \D_\Delta j_!\const{\R}{\Delta^\circ}
 \simeq Rj_*\const{\R}{\Delta^\circ}[n].
\]
The unit $\const{\R}{\Delta}\to
Rj_*\const{\R}{\Delta^\circ}$ is an isomorphism: the intersection of
each sufficiently small relative neighborhood with $\Delta^\circ$ is
nonempty and contractible, with no higher constant-sheaf cohomology.
Under these identifications the biduality map of
$\const{\R}{\Delta}$ is this unit.  Thus it is the canonical
biduality map, and is invertible, including at the boundary.  The case
$n=0$ is ordinary duality of the one-dimensional vector space.

Now take a simplex generator $A=Rf_*\const{\R}{\Delta}$.
Since $f$ is proper, \eqref{cat:verdier-formula} and
\eqref{cat:simplex-duality} give
\[
 \D_X A\simeq Rf_*\omega_\Delta.
\]
The boundary triangle \eqref{cat:simplex-boundary-triangle} expresses
$\omega_\Delta$ using constants on compact subpolyhedra of $\Delta$.
Consequently $\D_XA$ lies in the same generated category.

We check explicitly that proper duality transports the canonical
biduality map.  Put $F=Rf_*$ and $B=\const{\R}{\Delta}$, and write
\[
 \zeta_C:F\D_\Delta C\xrightarrow{\sim}\D_XFC
 \qquad(C=B,\D_\Delta B)
\]
for \eqref{cat:verdier-formula}.  The source calculation makes
$B,\D_\Delta B,\D_\Delta^2B$ bounded, as are their proper images.
The two comparisons $\zeta_B,\zeta_{\D_\Delta B}$ and the dual of
$\zeta_B$ therefore show that all terms in the following diagram are
bounded:
\begin{equation}
 \label{cat:proper-biduality}
 \begin{tikzcd}[column sep=large]
 FB \arrow[r,"\eta_{FB}"] \arrow[d,"F\eta_B"'] &
 \D_X^2FB \arrow[d,"\D_X(\zeta_B)"] \\
 F\D_\Delta^2B \arrow[r,"\zeta_{\D_\Delta B}"'] &
 \D_XF\D_\Delta B.
 \end{tikzcd}
\end{equation}
Under tensor--Hom adjunction, both composites correspond to the
pushed evaluation pairing with its two arguments interchanged.  The
cup map respects tensor symmetry, and both composites use the same
trace, so the diagram commutes.  Its left, bottom, and right arrows
are isomorphisms.  Hence $\eta_A=\eta_{FB}$ is an isomorphism.

The class of objects whose dual is bounded and belongs to the generated
category is closed under shifts and triangles.  On this class, the
objects for which $\eta$ is invertible are closed under these operations
as well: naturality gives a morphism of triangles, and an isomorphism on
two terms gives an isomorphism on the third.  The same assertions pass
to direct summands, by applying the dual functor to the splitting maps.
\Cref{cat:simplex-generators} now proves the duality assertions for all
four categories.

Each generator vanishes outside $f(\Delta)$.  Finite sums, triangles,
shifts, and summands preserve vanishing outside a finite union of such
compact sets.  Choose one such compact $K$ for $A$.  Open restriction of
duality shows that $\D_X A$ also vanishes on $X\setminus K$.
\end{proof}

\begin{remark}
\label{cat:restriction-caution}
There is no assertion that restriction preserves the test categories.
For example, $\const{\R}{[0,1]}$ belongs to $\T([0,1])$, whereas its
restriction to $(0,1)$ has no compact support there and hence does not
belong to $\T((0,1))$.  Restrictions in later germ arguments are taken
in the ambient derived categories.  Likewise, Verdier reflexivity in
\Cref{prop:test-duality} is not algebraic reflexivity of each stalk:
Verdier duality uses local compact-support information.  This distinction
allows proper continuous images with infinite-dimensional stalks.
\end{remark}

\begin{proposition}[Direct images and open extension]
\label{prop:test-pushforward}
For an allowed map $g:X\to Y$, direct image induces exact functors
\[
 g_*:\Tzero(X)\longrightarrow\Tzero(Y),
 \qquad g_*:\T(X)\longrightarrow\T(Y),
\]
and the analogous functors for $\Pzero,\Pcat$.  They are given by
$Rg_!\simeq Rg_*$ on these categories, commute coherently with Verdier
duality and its shifts, and are compatible with composition.
For an open inclusion $j:U\hookrightarrow X$, this functor is $j_!$ and
is fully faithful.
\end{proposition}

\begin{proof}
If $A$ vanishes outside a compact $K\subset X$, let $i:K\hookrightarrow
X$.  The canonical map $A\to i_*i^{-1}A$ is an isomorphism, as is checked
on cohomology stalks.  The map $g i:K\to Y$ is proper.  Composition of
direct images therefore identifies the natural comparison
\[
 Rg_!A\longrightarrow Rg_*A
\]
with the corresponding, invertible comparison for $g i$.  This argument
requires only compactness of $K$, not a polyhedral structure on $K$.
On a generator the image is
\[
 Rg_*Rf_*\const{\R}{B}\simeq R(gf)_*\const{\R}{B},
\]
which is again an allowed generator.  Exactness and preservation of
splittings prove category membership and boundedness for every object.
The semialgebraic case follows from the same composition check.

Both $A$ and $\D_XA$ have compact support by
\Cref{prop:test-duality}.  Hence \eqref{cat:verdier-formula} gives the
natural isomorphism
\begin{equation}
 \label{cat:push-duality}
 Rg_!\D_XA\xrightarrow{\sim}\D_Y Rg_!A.
\end{equation}
It preserves biduality because it is defined by evaluation and trace.
More explicitly, a pairing $\beta:A\otimes A\to\omega_X[-r]$ is sent
to the pairing
\begin{equation}
 \label{cat:push-pairing}
 \begin{split}
 Rg_!A\otimes Rg_!A
 &\longrightarrow Rg_!(A\otimes A)
 \xrightarrow{Rg_!\beta}Rg_!\omega_X[-r]
 \xrightarrow{\operatorname{Tr}_g[-r]}\omega_Y[-r].
 \end{split}
\end{equation}
The first arrow is the cup map obtained from the projection formula and
the adjunction map $g^{-1}Rg_!A\to A$ (using
$Rg_!A\to Rg_*A$).  Under tensor--Hom adjunction,
\eqref{cat:push-pairing} is exactly \eqref{cat:push-duality} composed
with the image of the adjoint of $\beta$.  The symmetry of tensor
products and the adjunction identities imply compatibility with
biduality.  For composable maps, the cup maps compose and
\[
 \operatorname{Tr}_{h g}
 =\operatorname{Tr}_h\circ Rh_!(\operatorname{Tr}_g)
\]
under the composition identifications.  Thus the duality functors and
their pairings compose coherently.  Notice that the last two arrows of
\eqref{cat:push-pairing} use $Rg_!$ on $\omega_X$; compact support of
$\omega_X$ is not required.

For an open inclusion, $Rj_!=j_!$ and $j^{-1}j_!$ is the identity.
The adjunction $(j_!,j^{-1})$ gives
\[
 \Hom_{\mathsf D(X)}(j_!A,j_!B)
 \simeq\Hom_{\mathsf D(U)}(A,B).
\]
Thus extension is fully faithful.  The already proved comparison with
$Rj_*$ applies because objects from the test category on $U$ have
compact support in $U$.
\end{proof}

\subsection{Shifted forms and the cone relation}

We use cochain shifts and homological Witt indexing:
\begin{equation}
 \label{cat:witt-indexing}
 \D_{X,r}=[-r]\D_X,
 \qquad W_r(\mathcal B)=W^{-r}(\mathcal B),
 \qquad E_r(X)=W_r(\T(X)).
\end{equation}
The biduality and exactness signs are the translated-duality signs of
triangular Witt theory; the exactness sign of $\D_{X,r}$ is
$\delta_r=(-1)^r$.  With these signs, a symmetric form of degree
$r$ is equivalently a pairing
\begin{equation}
 \label{cat:pairing-form}
 \beta:A\otimes A\longrightarrow\omega_X[-r],
 \qquad \beta\circ\tau_{A,A}=\beta,
\end{equation}
whose adjoint $\alpha:A\xrightarrow{\sim}\D_{X,r}A$ is invertible.
Here $\tau$ is tensor symmetry with the Koszul sign.  We give the
chain-level verification of these conventions in
\Cref{app:duality-signs}.  In particular, on an oriented $d$-manifold
the orientation identification $\omega_X[-d]\simeq\const{\R}{X}$
makes the multiplication pairing of the constant sheaf a symmetric
degree-$d$ form.  When its underlying object belongs to a test category,
it defines a class there.

We recall the one algebraic cone fact needed for homotopy invariance.
Write $\#$ for the chosen shifted duality.
Witt groups are generated under orthogonal sum by nonsingular symmetric
forms, with neutral forms set equal to zero; the negative of a pairing
is its additive inverse.  A triangular lagrangian for $(A,\alpha)$ is
an arrow $i:L\to A$ completed to a triangle
$L\to A\to\#L\to L[1]$ in which the second map is the composite
$A\xrightarrow{\alpha}\#A\xrightarrow{\#i}\#L$,
with the translated-duality signs, and the triangle is compatible with
its dual under the form and biduality, including the connecting
morphism. A form with such a lagrangian is neutral. The hyperbolic form
$H(L)$, on
$L\oplus\#L$, is an example.  These definitions and their equivalence
with the symmetric-cone description below are the triangular Witt
formalism of \cite{Balmer2000}; a convenient precise reference for the
cone statement is \cite[\S2.1, Lemma~2.1]{Woolf2008}.

Set $I=[0,1]$ and let $i_t:X\to I\times X$ be $x\mapsto(t,x)$,
for $t=0,1$.  For a degree-$r$ form $(A,\alpha)$ in $\T(X)$ or
$\Pcat(X)$, our goal is to prove
\[
 (i_0)_*[A,\alpha]=(i_1)_*[A,\alpha]
\]
in the corresponding Witt group of $I\times X$.
We will construct a neutral form on the sum of the two endpoint objects
whose diagonal pairings are a common nonzero real multiple of
$(i_0)_*\alpha$ and $-(i_1)_*\alpha$.  The symmetric-cone construction
below produces this form; pushing the resulting equality along a
homotopy will then identify the homomorphisms induced by its endpoints.

\begin{lemma}[Symmetric cone]
\label{cat:symmetric-cone}
Let $(\#,\eta)$ be a triangulated duality with exactness sign
$\delta\in\{1,-1\}$ on a category in which multiplication by $2$ is
invertible on every morphism group.  Let $u:L\to\#L$ be symmetric,
without an invertibility assumption.  In a triangle
\[
 L\xrightarrow{u}\#L\xrightarrow{v}C\xrightarrow{w}L[1],
\]
the cone has a nonsingular symmetric form
$\psi:C\xrightarrow{\sim}[1]\#C$ for the translated duality.
It is neutral and can be chosen to satisfy
\begin{equation}
 \label{cat:cone-compatibility}
 \psi v=-[1](\#w),
 \qquad [1](\#v)\psi=\delta[1](\eta_L)w,
\end{equation}
with the dual shift comparisons understood.  For
$\#=\D_{X,r+1}$ this is a neutral form of degree $r$.
\end{lemma}

\begin{proof}
Dualize and rotate the triangle.  Its first arrow is $\delta\#u$.
Symmetry of $u$ makes its first two terms compatible with
$\delta\eta_L$ and the identity of $\#L$, with fourth vertical map
$\delta[1](\eta_L)$.  Complete these maps to a morphism of triangles
with third map $\psi_0$.
The signed transpose of $\psi_0$ is another such completion.
Replacing $\psi_0$ by
\[
 \frac12\bigl(\psi_0-\delta[1](\#\psi_0)\eta_C\bigr)
\]
preserves the first two squares and imposes symmetry for the translated
biduality $-\delta\eta$.  The resulting map $\psi$ satisfies
\eqref{cat:cone-compatibility}.  It is an isomorphism by the
two-out-of-three property for the morphism of triangles.  By definition
it is the symmetric boundary of $u$, hence neutral.  Equivalently, rotate
the displayed triangle to start with $\#L\to C$; its third object is
$L[1]\simeq[1]\#(\#L)$, and \eqref{cat:cone-compatibility} gives the
triangular lagrangian.
All these operations take place in the given full triangulated
category; averaging uses that $2$ is invertible.  Finally
$[1]\D_{X,r+1}=\D_{X,r}$ with exactly the translated signs, proving
the degree assertion.
\end{proof}

\subsection{Natural external products}

To compare the two endpoint copies of a form on $X$, we will combine
it with the interval boundary triangle.  The next lemma verifies that
these products remain in the test category and carry the required
pairings.

Let $e_t:\{t\}\hookrightarrow I$ for $t=0,1$.
Write
\begin{equation}
 \label{cat:interval-objects}
 H=\const{\R}{I},\qquad
 J=j_!\const{\R}{(0,1)},\qquad
 E_t=(e_t)_*\R,\qquad E=E_0\oplus E_1,
\end{equation}
where $j:(0,1)\hookrightarrow I$.  Let $\mathcal C_I$ be the full
thick subcategory generated by $H,E_0,E_1$.  It is contained in
$\Pcat(I)$ and $\T(I)$, and contains $J$ by the triangle
\begin{equation}
 \label{cat:interval-triangle}
 J\xrightarrow{u}H\xrightarrow{v}E
 \xrightarrow{\partial}J[1].
\end{equation}

\begin{lemma}[Products with the interval]
\label{cat:interval-products}
For $A\in\T(X)$ and $B\in\mathcal C_I$, the external product
$B\boxtimes A$ belongs to $\T(I\times X)$; the same assertion holds
for $\Pcat$.  For all integers $s,r$, the natural map
\begin{equation}
 \label{cat:external-duality}
 \mu_{s,r}:\D_{I,s}B\boxtimes\D_{X,r}A
 \xrightarrow{\sim}
 \D_{I\times X,s+r}(B\boxtimes A)
\end{equation}
is an isomorphism.  The map is adjoint to the product of evaluations,
with the tensor shift signs; it commutes with proper direct image,
duality, and the trace maps.  In particular the product of symmetric
pairings is symmetric for the sum of their degrees.
\end{lemma}

\begin{proof}
First suppose $A=Rf_*\const{\R}{\Delta}$.  For $B=H$ the
proper external-image formula gives
\[
 H\boxtimes A
 \simeq R(1\times f)_*\const{\R}{I\times\Delta}.
\]
The source is a compact finite polyhedron.  For $B=E_t$ the product is
the generator on the endpoint copy of $\Delta$.  The triangle
\eqref{cat:interval-triangle} proves membership for $B=J$ as well.
Exactness and preservation of summands in each variable prove the
membership statement for all $B,A$.  All source maps in this argument
are semialgebraic when $f$ is.

We describe the natural map in \eqref{cat:external-duality} before
proving that it is invertible.  For locally compact spaces $Z,X$ there
is a product trace map
\begin{equation}
 \label{cat:product-trace}
 t_{Z,X}:\omega_Z\boxtimes\omega_X\longrightarrow\omega_{Z\times X}.
\end{equation}
By the adjunction for $a_{Z\times X}$, this is the map corresponding,
under compact-support K\"unneth, to the tensor of the two traces to
$\R$.  To define $\mu_{s,r}$, pair
\[
 (\D_{Z,s}B\boxtimes\D_{X,r}A)
       \otimes(B\boxtimes A)
\]
by moving the second dual factor past $B$, evaluating the two pairs,
combining their output shifts, and applying $t_{Z,X}[-s-r]$.
Take its tensor--Hom adjoint.  This construction proves naturality in
both variables and specifies the actual comparison, rather than merely
an abstract isomorphism of its objects.

The construction is compatible with proper pushforward.  We verify the
two ingredients of this assertion.  First, for a proper $f:Z'\to Z$,
the two composites from
$R(f\times1)_!(\omega_{Z'}\boxtimes\omega_X)$ to
$\omega_{Z\times X}$ agree: one uses $\operatorname{Tr}_f$ and then
the product trace; the other uses the source product trace and then
$\operatorname{Tr}_{f\times1}$.  Under adjunction to the point and
K\"unneth, both are
\[
 \operatorname{Tr}_{a_{Z'}}\otimes\operatorname{Tr}_{a_X}.
\]
Their equality therefore follows from adjunction, including its
composition identities.  Second, the cup maps used to push evaluation
in \eqref{cat:push-pairing} commute with external products.  Indeed
the external-image isomorphism is obtained from proper base change and
the projection formula, and after inverse-image adjunction its cup map
is the tensor of the two counits.  Tensoring these counits before or
after evaluating gives the same map.  These two equalities show that
the natural $\mu$ is carried to the natural $\mu$ by proper-image
duality.  They also show its compatibility with biduality.  The
assertion for shifts follows by inserting the tensor shift
comparisons, not by changing any trace.

It remains to prove invertibility on the objects in the lemma. First
take $s=r=0$; the specified shift comparisons give the other degrees. On
$I\times\Delta$, for $B=H$ and $A=\const{\R}{\Delta}$,
\eqref{cat:simplex-duality} identifies both sides with the constant
sheaf on $(0,1)\times\Delta^\circ$ extended by zero and shifted by
$1+\dim\Delta$.  This calculation applies to the product polyhedron
with its boundary, using its product orientation.  The map on the
interior is the product of the orientation evaluations and is an
isomorphism; both sides vanish on the boundary.  For $B=E_t$, closed
direct-image duality reduces the map to the identity duality comparison
on $\Delta$.  Thus it is an isomorphism for $H,E_0,E_1$.
Triangles and summands prove the assertion for every
$B\in\mathcal C_I$ and for their shifts.

For $A=Rf_*\const{\R}{\Delta}$, the external-image isomorphism
and proper duality identify \eqref{cat:external-duality} with the
proper image of the comparison just verified on $I\times\Delta$.
The preceding naturality check is what identifies it with that map.
It is consequently an isomorphism even for an arbitrary continuous
$f$.  Finally, for fixed $B$ the locus of $A$ for which this natural
comparison is an isomorphism is closed under triangles, shifts, and
summands.  \Cref{cat:simplex-generators} completes the proof.

For the last assertion, write the product pairing as the tensor of its
two pairings, preceded by the middle-factor permutation and followed
by \eqref{cat:product-trace}.  Naturality of the Koszul symmetry shows
that interchanging its two arguments is the same as interchanging the
arguments of each factor pairing.  Both factor pairings are symmetric,
so the product is symmetric.
\end{proof}

\subsection{The interval boundary and homotopies}

Orient $I$ from $0$ to $1$.  The local calculation already used gives
\begin{equation}
 \label{cat:interval-duals}
 \D_{I,1}J\simeq H,\qquad
 \D_{I,1}H\simeq J,\qquad
 \D_I E\simeq E.
\end{equation}
The first arrow $u$ of \eqref{cat:interval-triangle}, considered as
$J\to\D_{I,1}J$, is symmetric.  Its pairing is multiplication
$J\otimes J\to J=\omega_I[-1]$.

\begin{lemma}[The two endpoint forms]
\label{cat:endpoint-forms}
Let $(A,\alpha)$ be a symmetric degree-$r$ form in $\T(X)$ or
$\Pcat(X)$, and let $i_t:X\to I\times X$ be $x\mapsto(t,x)$.
There is a neutral degree-$r$ form on
\[
 E\boxtimes A\simeq(i_0)_*A\oplus(i_1)_*A
\]
whose two diagonal forms are $c(i_0)_*\alpha$ and
$-c(i_1)_*\alpha$, for a common nonzero real scalar $c$.  There are no
off-diagonal terms.
\end{lemma}

\begin{proof}
We first determine the interval boundary itself.  Compact-support
cohomology of \eqref{cat:interval-triangle} is the exact sequence
\begin{equation}
 \label{cat:interval-cohomology}
 0\longrightarrow\R
 \xrightarrow{x\mapsto(x,x)}\R^2
 \xrightarrow{(x_0,x_1)\mapsto x_1-x_0}\R
 \longrightarrow0,
\end{equation}
where the last identification uses the chosen orientation of $I$.
For example, it follows by taking the cellular cochains of the pair
$(I,\{0,1\})$: the coboundary of endpoint values is their difference.
Proper duality to the point therefore identifies the shifted dual of
$\partial$ with a map $H\to E$ having opposite nonzero coefficients
at the two ends.  This determines the sheaf morphism, since
\[
 \Hom_{\mathsf D(I)}(H,E_t)
 \simeq\Hom_{\mathsf D(\pt)}(\R,\R)=\R
\]
by closed direct-image adjunction, and taking global sections detects
these two coefficients.  The cone identity
\eqref{cat:cone-compatibility} consequently gives opposite endpoint
pairings for the symmetric cone of $u$.

Now use the symmetric morphism supplied by the product of $u$ and
$\alpha$:
\begin{equation}
 \label{cat:cylinder-morphism}
 J\boxtimes A
 \xrightarrow{u\boxtimes\alpha}
 H\boxtimes\D_{X,r}A
 \xrightarrow{\mu_{1,r}}
 \D_{I\times X,r+1}(J\boxtimes A).
\end{equation}
It lies in the required category by
\Cref{cat:interval-products}, and it is symmetric by that lemma.
Identifying its middle object with $H\boxtimes A$ using $\alpha$,
its cone triangle is
\[
 J\boxtimes A\longrightarrow H\boxtimes A
 \xrightarrow{v\boxtimes1}E\boxtimes A
 \xrightarrow{\partial\boxtimes1}(J\boxtimes A)[1].
\]
\Cref{cat:symmetric-cone} supplies a neutral nonsingular form $\psi$
of degree $r$ on $E\boxtimes A$.

The two slices are disjoint closed subsets.  Closed direct-image
adjunction thus makes every morphism between their supported
complexes zero, in every degree.  In particular $\psi$ is diagonal.
On either diagonal its precomposition with $v\boxtimes1$ determines
it: restriction of $H\boxtimes A$ to that slice is $A$, and the
restriction of $v\boxtimes1$ is the identity.  Apply
\eqref{cat:cone-compatibility} to compute this precomposition.
By the natural product comparison, the shifted dual of
$\partial\boxtimes1$ is the product of the shifted dual of
$\partial$ with the identity of $\D_{X,r}A$, up to a single shift
sign independent of the endpoint, as calculated in
\Cref{cat:product-shift-sign}.
After precomposing with $1\boxtimes\alpha$, the two coefficients
are therefore a common nonzero scalar times $\alpha$ and its
negative, as asserted.  \Cref{app:duality-signs} fixes the absolute
normalization of this scalar.
\end{proof}

\begin{theorem}[Homotopy invariance]
\label{thm:homotopy}
If $h:I\times X\to Y$ is a continuous homotopy of maps of test
spaces, its endpoint maps induce equal homomorphisms
\[
 h_{0*}=h_{1*}:E_r(X)\longrightarrow E_r(Y)
 \qquad(r\in\Z).
\]
The same assertion holds for $W_r(\Pcat(-))$ under the allowed
semialgebraic homotopies.  In particular, homotopy equivalences induce
isomorphisms, and a nonempty contractible test space has the Witt groups
of a point.  The semialgebraic assertion uses a contraction in the
allowed semialgebraic category.
\end{theorem}

\begin{proof}
For a form $(A,\alpha)$, \Cref{cat:endpoint-forms} gives
\[
 [\,(i_0)_*A,c(i_0)_*\alpha\,]
 +[\,(i_1)_*A,-c(i_1)_*\alpha\,]=0
\]
in the degree-$r$ Witt group of the cylinder.  Multiplying a pairing
by a positive real scalar does not change its isometry class: scalar
multiplication of its underlying object by the inverse square root
gives an isometry.  If $c<0$, the displayed relation is the negative
of the one with $|c|$.  Thus in either case
\[
 (i_0)_*[A,\alpha]=(i_1)_*[A,\alpha].
\]
Push this equality along $h$ using
\Cref{prop:test-pushforward}.  Its composite maps are $h_0$ and
$h_1$, so it proves the desired equality on form classes and hence on
the Witt group.  Properness of the whole homotopy is unnecessary:
every object in the relation has support in $I\times K$ for a compact
$K\subset X$.  The proof for $\Pcat$ uses exactly the same relation,
whose compact-source maps and homotopy images are then semialgebraic.
Apply the endpoint equality to the two compositions of a homotopy
equivalence and its inverse to obtain the final assertions.
\end{proof}

This proof implements the interval argument of
\cite[Proposition~3.11]{Woolf2008} inside the actual generated
categories.  The continuous maps in their generators have only been
used through proper direct image, so no constructibility condition on
their images was introduced.

At a point both test categories are exactly the bounded derived category
of finite-dimensional real vector spaces: a compact polyhedron has
finite-dimensional bounded cohomology, and the point generator and its
finite sums and shifts already generate every such complex.  We will
use this identification and its elementary Witt calculation in
\Cref{sec:lattice}.
